\documentclass[11pt,a4paper]{article}
\usepackage[utf8]{inputenc}
\usepackage[top=0.9in, bottom=0.8in, left=1in, right=1in]{geometry}
\usepackage{hyperref}

\setlength{\parskip}{0.8em}
\setlength{\parindent}{0pt}

\begin{document}

\begin{center}
    {\LARGE \textbf{Liuxi Mei}} \\
    \vspace{2mm}
    Linköping, Sweden $\cdot$ liume102@student.liu.se $\cdot$ +46-0704898962
    \vspace{4mm}
\end{center}

\textbf{Date:} April 5, 2026 \\
\textbf{To:} The Hiring Committee \\
\textbf{Subject:} Application for Doctoral Researcher in AI \& Bioinformatics \\

Dear Hiring Committee,

Please accept my application for the Doctoral Researcher position in AI and Bioinformatics. I am currently completing my master's degree and will officially graduate and receive my diploma in June of this year.

Currently, the combination of AI and biology is pushing science forward. Biology always requires clear explanations, human review, and lab validation. However, reliable AI tools can deeply speed up these early steps. Huge amounts of multi-omics data exist today with massive unused potential. Digging into this complex data can reveal many hidden medical discoveries.

My background brings a unique advantage: over five years of software industry experience combined with recent academic research. In the software field, managing massive real-time stream data was my daily focus. These strong engineering habits proved very useful in my master's studies. For my thesis, processing complex biological datasets became my main task, specifically using LINCS L1000 gene expression data and OmniPath networks.

Building on this data, my academic research focused heavily on Graph Neural Networks (GNNs). Industry work previously taught me how to put stable machine learning models into real production. This practical mindset shaped my thesis on modeling high-throughput experimental data. By analyzing the GNN outputs, the project successfully mapped and explained biological pathways. Seeing these clear biological results from mathematical models was amazing. This is exactly why I chose to apply for your specific PhD topic. I want to take this graph-based pathway research much further. 

Ultimately, my professional ambition is to bridge the gap between complex AI systems and real medical discoveries. My software background proves I can handle the technical scale, and my recent thesis drives my passion for bioinformatics. Your project's academic focus aligns perfectly with this goal. Thank you for reading my application. I look forward to speaking with you soon.

Sincerely,

Liuxi Mei

\vspace{0.5cm}
\noindent
\begin{minipage}{\textwidth}
\textbf{References:}

\vspace{0.3cm}
\begin{minipage}[t]{0.48\textwidth}
\textbf{Oleg Sysoev}\\
Senior Associate Professor, Linköping University\\
Department of Computer and Information Science (IDA)\\
Master's Thesis Supervisor\\
Email: \href{mailto:oleg.sysoev@liu.se}{oleg.sysoev@liu.se}
\end{minipage}%
\hfill
\begin{minipage}[t]{0.48\textwidth}
\textbf{Krzysztof Bartoszek}\\
Senior Associate Professor, Linköping University\\
Department of Computer and Information Science (IDA)\\
Bioinformatics Teacher\\
Email: \href{mailto:krzysztof.bartoszek@liu.se}{krzysztof.bartoszek@liu.se}
\end{minipage}
\end{minipage}

\end{document}